% Compile from this directory:
%   pdflatex financial_sde_report.tex
%   pdflatex financial_sde_report.tex
% The experiment artifacts are regenerated from the repository root with:
%   python financial_sde_project.py
\documentclass[11pt,a4paper]{article}

\usepackage[utf8]{inputenc}
\usepackage[T1]{fontenc}
\usepackage{amsmath,amssymb,amsthm,bm,mathtools}
\usepackage[a4paper,margin=2.35cm]{geometry}
\usepackage{graphicx}
\graphicspath{{generated/}}
\usepackage{booktabs,array,longtable,multirow}
\usepackage{listings}
\usepackage{xcolor}
\usepackage{hyperref}
\usepackage{enumitem}
\usepackage{float}
\usepackage{caption}
\usepackage{subcaption}
\usepackage{algorithm}
\usepackage{algpseudocode}
\usepackage{mdframed}
\usepackage{microtype}
\usepackage{mathptmx}
\usepackage{url}

\definecolor{bluecb}{RGB}{0,119,187}
\definecolor{orangecb}{RGB}{238,119,51}
\definecolor{formbg}{RGB}{250,250,245}
\definecolor{formframe}{RGB}{180,170,140}
\hypersetup{colorlinks=true,linkcolor=bluecb,citecolor=bluecb,urlcolor=bluecb}
\setlength{\parindent}{1.5em}
\setlength{\parskip}{0.25em}
\captionsetup{font=small,labelfont=bf}

\lstset{
  language=Python,
  basicstyle=\ttfamily\scriptsize,
  keywordstyle=\color{blue}\bfseries,
  commentstyle=\color{green!40!black},
  stringstyle=\color{red!65!black},
  numbers=left,
  numberstyle=\tiny\color{gray},
  frame=single,
  breaklines=true,
  showstringspaces=false,
  tabsize=4,
  captionpos=b
}

\newtheorem{theorem}{Theorem}[section]
\newtheorem{proposition}[theorem]{Proposition}
\theoremstyle{definition}
\newtheorem{definition}[theorem]{Definition}
\newtheorem{assumption}[theorem]{Assumption}
\theoremstyle{remark}
\newtheorem{remark}[theorem]{Remark}

\newmdenv[
  backgroundcolor=formbg,linecolor=formframe,linewidth=0.8pt,
  roundcorner=3pt,innermargin=6pt,outermargin=5pt,
  innertopmargin=7pt,innerbottommargin=7pt
]{formbox}

\algrenewcommand{\algorithmicrequire}{\textbf{Input:}}
\algrenewcommand{\algorithmicensure}{\textbf{Output:}}

\title{\textbf{Reliable Numerical Simulation of Financial SDEs:}\\
\textbf{GBM Benchmarks and a Non-Affine Volatility Model}\\[2mm]
\large Project 1 (ODE IVP Extension Pathway: Financial SDE)}
\author{
  \texttt{Team H}\\
  \small \texttt{Jinyang Wang, 2364059 (Project Manager)} \and
  \texttt{Yangxiaoxiang Zheng, 1929725 (Mathematical Theory)}\\
  \small \texttt{Jinsen Chen, 2363993 (Algorithm Implementation)} \and
  \texttt{Zheyuan Zheng, 2362114 (Visualization \& Report)}\\
  \small \texttt{Dongdong Hao, 2364312 (Testing \& Validation)}\\
  \small Department of Applied Mathematics
}
\date{18 September 2026}

\begin{document}
\maketitle

\begin{abstract}
This project develops and validates a reproducible numerical workflow for two
financial stochastic differential equations (SDEs).  Geometric Brownian motion
(GBM) is first used as an exact benchmark: Euler--Maruyama (EM), scalar Milstein,
and exact lognormal transitions are driven by identical Brownian increments.
This coupling isolates time-discretisation error from sampling noise.  With
$300{,}000$ paths, the fitted terminal strong-convergence slopes are
$0.4899\pm0.0026$ for EM and $0.9099\pm0.0167$ for Milstein, consistent with the
theoretical orders $1/2$ and $1$.  The exactly evaluable bias in
$\mathbb{E}[S_T]$ has fitted slope $0.9982\pm0.0004$.  At a matched mean absolute
error target of $0.85$ price units, EM uses 128 path steps while Milstein uses
2 under the declared step-count cost measure.  The experiment also checks the
terminal lognormal law and demonstrates that price-space updates are not
unconditionally positivity preserving.  We then solve a non-affine
stochastic-volatility model with bounded volatility in log price.  Coarse paths are coupled to an
$N_{\mathrm{ref}}=2048$ EM reference by summing fine-grid Brownian increments.
The empirical strong slope is $0.5111\pm0.0047$; Brownian variance and
correlation diagnostics pass, and all checked reference terminal states are
finite, with positive terminal prices.  Because the reference is numerical and call-payoff bias
confidence intervals include zero, no unsupported weak-order claim is made.
\end{abstract}

\noindent\textbf{Keywords:} stochastic differential equations; geometric
Brownian motion; Euler--Maruyama; Milstein method; strong convergence; weak
convergence; stochastic volatility; Monte Carlo coupling

\tableofcontents
\newpage

% ============================================================
\section{Introduction and motivation}
% ============================================================

\subsection{Problem background}
An SDE model is useful only if its numerical error can be separated from Monte
Carlo uncertainty.  GBM is therefore more than a toy model: its exact pathwise
transition, moments, and terminal distribution provide an oracle against which
an implementation can be tested.  Higham's algorithmic treatment emphasises
precisely this role of coupled paths when illustrating EM and Milstein
\cite{higham2001}; the broader convergence theory is developed by Kloeden and
Platen \cite{kloedenplaten1992} and Milstein \cite{milstein1995}.

The project asks a second, more practically relevant question: what changes
when no exact transition is available?  We answer it with the prescribed
two-factor non-affine stochastic-volatility model.  The same Brownian path is
represented on nested grids, so agreement under refinement becomes evidence
for convergence.  It is not relabelled as an exact error.  This distinction is
central to reliable computational finance and underlies multilevel coupling
methods \cite{giles2008}.

\subsection{Research questions and contribution}
The numerical investigation addresses four questions.
\begin{enumerate}[leftmargin=2.2em]
  \item Do coupled simulations recover the predicted GBM strong orders
  $1/2$ (EM) and $1$ (Milstein), and weak order $1$ for the mean?
  \item At a matched strong-error target, which method uses fewer time-step
  coefficient evaluations?
  \item When can price-space updates lose positivity, and does the simulated
  terminal law agree with the exact lognormal distribution?
  \item For the non-affine model, do nested-grid EM solutions converge, pass
  Brownian-increment checks, remain finite and positive, and use a sufficiently
  resolved numerical reference?
\end{enumerate}

The main contribution is a single, auditable script that implements every
scheme, fixes all seeds, regenerates every figure and printed number, and runs
explicit PASS/FAIL checks.  The analysis deliberately separates theorem,
measurement, and limitation.

\subsection{Report map}
Section~\ref{sec:model} defines both models.  Section~\ref{sec:methods} derives
the numerical schemes and convergence estimators.  Section~\ref{sec:implementation}
documents the implementation and validation design.  Section~\ref{sec:results}
presents GBM and non-affine experiments.  Section~\ref{sec:limitations}
identifies what the evidence does not establish, and Section~\ref{sec:conclusion}
concludes.  The complete source and AI/individual-contribution records appear
in the appendices.

% ============================================================
\section{Financial SDE models}\label{sec:model}
% ============================================================

\subsection{Geometric Brownian motion}
Let $(W_t)_{t\ge0}$ be a standard Brownian motion.  Under the physical measure,
GBM is
\begin{equation}
  dS_t=\mu S_t\,dt+\sigma S_t\,dW_t,\qquad S_0>0.
  \label{eq:gbm}
\end{equation}
For $f(s)=\log s$, It\^o's formula gives
\begin{align}
  d\log S_t
  &=f'(S_t)dS_t+\tfrac12 f''(S_t)(dS_t)^2 \notag\\
  &=\left(\mu-\tfrac12\sigma^2\right)dt+\sigma dW_t.
  \label{eq:ito-log}
\end{align}
Integrating~\eqref{eq:ito-log} yields the exact solution
\begin{equation}
  S_t=S_0\exp\!\left[\left(\mu-\tfrac12\sigma^2\right)t+\sigma W_t\right].
  \label{eq:gbm-exact}
\end{equation}
Consequently,
\begin{equation}
  \mathbb E[S_T]=S_0e^{\mu T},\qquad
  \operatorname{Var}(S_T)=S_0^2e^{2\mu T}\bigl(e^{\sigma^2T}-1\bigr),
  \label{eq:gbm-moments}
\end{equation}
and $S_T$ is lognormal.  The reproducible benchmark uses
\begin{equation}
  S_0=100,\qquad \mu=0.05,\qquad \sigma=0.40,\qquad T=1.
  \label{eq:gbm-params}
\end{equation}
The volatility is intentionally high enough to make the coarse-step positivity
question visible while keeping a familiar annual horizon.

\subsection{Non-affine stochastic volatility}
Write $X_t=\log S_t$.  The prescribed model is
\begin{align}
 dX_t &= \left(\mu-\tfrac12g(Y_t)^2\right)dt+g(Y_t)dW_t^{(1)},
 \label{eq:nonaff-x}\\
 dY_t &= \kappa(\theta-Y_t)dt+\xi\sqrt{1+Y_t^2}\,dW_t^{(2)},
 \label{eq:nonaff-y}\\
 g(y)&=\sigma_{\min}+\frac{\sigma_{\max}-\sigma_{\min}}{1+e^{-y}},
 \qquad d\langle W^{(1)},W^{(2)}\rangle_t=\rho\,dt.
 \label{eq:logistic}
\end{align}
The benchmark parameters are
\begin{equation}
\begin{gathered}
S_0=100,\quad Y_0=0,\quad \mu=0.05,\quad \kappa=2,
\quad\theta=-0.2,\quad\xi=0.6,\quad\rho=-0.7,\\
\sigma_{\min}=0.10,\quad\sigma_{\max}=0.50,\quad T=1.
\end{gathered}
\label{eq:nonaff-params}
\end{equation}
The logistic map guarantees $0.10<g(y)<0.50$.  Its derivative is bounded,
and $y\mapsto\sqrt{1+y^2}$ is globally Lipschitz because the magnitude of its
derivative is at most one.  Thus the coefficients have the regularity and
linear-growth structure needed for standard EM strong-convergence arguments;
nevertheless, the finite reference used below remains a numerical benchmark,
not a closed-form oracle.

% ============================================================
\section{Discretisation and mathematical analysis}\label{sec:methods}
% ============================================================

\subsection{Euler--Maruyama for GBM}
On $t_n=nh$, $h=T/N$, let $\Delta W_n=W_{t_{n+1}}-W_{t_n}\sim
\mathcal N(0,h)$.  Freezing both coefficients at $S_n$ and integrating over one
step gives
\begin{equation}
  S_{n+1}^{\mathrm{EM}}
  =S_n^{\mathrm{EM}}+\mu S_n^{\mathrm{EM}}h
   +\sigma S_n^{\mathrm{EM}}\Delta W_n
  =S_n^{\mathrm{EM}}(1+\mu h+\sigma\Delta W_n).
  \label{eq:em-gbm}
\end{equation}
For globally Lipschitz drift and diffusion, EM has strong order $1/2$ and weak
order $1$ for sufficiently smooth test functions \cite{higham2001,kloedenplaten1992}.
GBM coefficients are globally Lipschitz, so the standard result applies.

\subsection{Milstein for GBM}
The scalar It\^o--Taylor expansion retains the iterated stochastic integral
\begin{equation}
 \int_{t_n}^{t_{n+1}}\int_{t_n}^{s}dW_r\,dW_s
 =\tfrac12\left[(\Delta W_n)^2-h\right].
 \label{eq:ito-iterated}
\end{equation}
For diffusion $b(s)=\sigma s$, $b'(s)=\sigma$, so the Milstein update is
\begin{equation}
 S_{n+1}^{\mathrm{Mil}}
 =S_n^{\mathrm{Mil}}+\mu S_n^{\mathrm{Mil}}h
 +\sigma S_n^{\mathrm{Mil}}\Delta W_n
 +\tfrac12\sigma^2S_n^{\mathrm{Mil}}\bigl[(\Delta W_n)^2-h\bigr].
 \label{eq:milstein-gbm}
\end{equation}
Including~\eqref{eq:ito-iterated} cancels the leading stochastic Taylor term
omitted by EM, raising the scalar strong order to $1$ while retaining weak order
$1$ \cite{milstein1995}.

\subsection{Exact transition and pathwise coupling}
For the same increment used in~\eqref{eq:em-gbm}--\eqref{eq:milstein-gbm},
\begin{equation}
 S_{n+1}^{\mathrm{exact}}
 =S_n^{\mathrm{exact}}
 \exp\!\left[\left(\mu-\tfrac12\sigma^2\right)h+\sigma\Delta W_n\right].
 \label{eq:exact-step}
\end{equation}
The exact, EM, and Milstein paths therefore share all $\Delta W_n$.  The terminal
strong-error estimator is
\begin{equation}
 \widehat e_{\mathrm{s}}(h)
 =\frac1M\sum_{m=1}^{M}\left|S_{N,m}^{(h)}-S_{T,m}^{\mathrm{exact}}\right|,
 \label{eq:strong-estimator}
\end{equation}
and a log--log least-squares slope estimates $p$ in
$\mathbb E|S_N-S_T|=O(h^p)$.  Pairing is essential: independent exact paths
would add $O(M^{-1/2})$ sampling variability unrelated to discretisation.

\subsection{Weak error and an analytic bias oracle}
For a test function $\varphi$, weak error is
\begin{equation}
 e_{\mathrm{w}}(h)=\left|\mathbb E\varphi(S_N^{(h)})-
 \mathbb E\varphi(S_T)\right|.
 \label{eq:weak-def}
\end{equation}
For $\varphi(s)=s$, both EM and Milstein step multipliers have conditional mean
$1+\mu h$.  Hence
\begin{equation}
 \mathbb E[S_N^{\mathrm{EM}}]
 =\mathbb E[S_N^{\mathrm{Mil}}]
 =S_0(1+\mu h)^N.
 \label{eq:numerical-mean}
\end{equation}
Using $N=T/h$ and $\log(1+\mu h)=\mu h-\tfrac12\mu^2h^2+O(h^3)$,
\begin{align}
 S_0(1+\mu h)^{T/h}
 &=S_0\exp\!\left(\mu T-\tfrac12\mu^2Th+O(h^2)\right),\\
 e_{\mathrm w}(h)
 &=\tfrac12S_0e^{\mu T}\mu^2T\,h+O(h^2).
 \label{eq:weak-expansion}
\end{align}
Thus the mean bias has exact asymptotic order one.  The code reports both this
deterministic bias and a paired Monte Carlo confidence interval for
$S_N-S_T$; the former prevents sampling noise from hiding the weak rate.

\subsection{Positivity analysis}
EM is positive for a step only when
\begin{equation}
 1+\mu h+\sigma\sqrt h Z_n>0,
\end{equation}
so the one-step failure probability is
\begin{equation}
 p_{\mathrm{neg,EM}}(h)
 =\Phi\!\left(-\frac{1+\mu h}{\sigma\sqrt h}\right)>0.
 \label{eq:em-neg-prob}
\end{equation}
For Milstein, the multiplier is the quadratic
\begin{equation}
 q(z)=\tfrac12\sigma^2hz^2+\sigma\sqrt h z
 +1+\mu h-\tfrac12\sigma^2h.
 \label{eq:milstein-quadratic}
\end{equation}
Its minimum is
\begin{equation}
 q_{\min}=\tfrac12+\mu h-\tfrac12\sigma^2h.
 \label{eq:milstein-min}
\end{equation}
For~\eqref{eq:gbm-params} and $0<h\le1$, this is positive, explaining why no
Milstein positivity violation is observed here.  Equation~\eqref{eq:milstein-min}
also shows that this is parameter- and step-dependent, not unconditional
positivity preservation.  Exact/log-price simulation remains positive by
construction.

\subsection{Correlated EM for the non-affine model}
Generate independent $Z_{1,n},Z_{2,n}\sim\mathcal N(0,1)$ and set
\begin{align}
 \Delta W_n^{(1)}&=\sqrt h Z_{1,n},\\
 \Delta W_n^{(2)}&=\sqrt h\left(\rho Z_{1,n}+\sqrt{1-\rho^2}Z_{2,n}\right).
 \label{eq:corr-inc}
\end{align}
The EM update is
\begin{align}
 X_{n+1}&=X_n+\left(\mu-\tfrac12g(Y_n)^2\right)h
              +g(Y_n)\Delta W_n^{(1)},
 \label{eq:em-x}\\
 Y_{n+1}&=Y_n+\kappa(\theta-Y_n)h
              +\xi\sqrt{1+Y_n^2}\Delta W_n^{(2)},
 \label{eq:em-y}\\
 S_{n+1}&=e^{X_{n+1}}>0.
 \label{eq:recover-s}
\end{align}
The logistic function is evaluated with two algebraically equivalent branches
to avoid overflow when $|Y_n|$ is large.

A scalar Milstein formula is \emph{not} applied to
\eqref{eq:nonaff-x}--\eqref{eq:nonaff-y}.  Written in two independent Brownian
drivers, the diffusion vector fields have nonzero cross derivatives through
$g(Y)$ and $\sqrt{1+Y^2}$.  A full multidimensional Milstein discretisation
would generally require cross terms and simulated iterated stochastic
integrals.  Omitting them while claiming strong order one would be
mathematically unjustified.

\subsection{Nested-grid numerical reference}
Let $h_{\mathrm{ref}}=T/N_{\mathrm{ref}}$.  For a coarse step $h=Kh_{\mathrm{ref}}$,
the coupling is
\begin{equation}
 \Delta W_{n,\mathrm{coarse}}^{(i)}
 =\sum_{j=0}^{K-1}\Delta W_{Kn+j,\mathrm{ref}}^{(i)},\qquad i=1,2.
 \label{eq:nested-sum}
\end{equation}
Therefore coarse and fine solutions represent the same Brownian path.  We use
$N_{\mathrm{ref}}=2048$ for the main experiment and independently compare it
with 4096 and 8192 steps on 4000 fresh, mutually coupled paths.
The same validation sample is used to assess sensitivity to the reference grid.  For the call payoff
$P(S_T)=(S_T-100)^+$, paired differences estimate weak bias and substantially
reduce variance.  We report a weak trend only where its confidence interval is
separated from zero.

\subsection{Uncertainty quantification}
For independent path-level observations $D_1,\ldots,D_M$, every displayed
Monte Carlo mean uses
\begin{equation}
 \overline D\ \pm\ 1.96\frac{s_D}{\sqrt M},
 \qquad
 s_D^2=\frac{1}{M-1}\sum_{m=1}^{M}(D_m-\overline D)^2.
 \label{eq:ci}
\end{equation}
The bands around fitted log--log lines are regression diagnostics, not repeated
Monte Carlo uncertainty.  Their slope standard errors quantify linear-fit
uncertainty conditional on the chosen mesh range.

% ============================================================
\section{Implementation and reproducibility}\label{sec:implementation}
% ============================================================

\subsection{Core algorithm}
Algorithm~\ref{alg:nested} states the non-affine nested-grid calculation.  GBM
uses the same coupling pattern with the exact step~\eqref{eq:exact-step}.

\begin{algorithm}[H]
\caption{Coupled non-affine EM convergence experiment}
\label{alg:nested}
\begin{algorithmic}[1]
\Require Parameters~\eqref{eq:nonaff-params}, seed, path count $M$, fine count
$N_{\rm ref}$, coarse counts $\mathcal N$
\Ensure Strong errors, paired payoff biases, confidence intervals, diagnostics
\For{each memory-safe batch of paths}
  \State draw independent fine-grid arrays $Z_1,Z_2$
  \State construct correlated increments using~\eqref{eq:corr-inc}
  \State integrate~\eqref{eq:em-x}--\eqref{eq:em-y} on the reference grid
  \For{$N\in\mathcal N$}
    \State sum groups of $N_{\rm ref}/N$ increments using~\eqref{eq:nested-sum}
    \State integrate the coarse path and store paired terminal differences
  \EndFor
\EndFor
\State estimate means and confidence intervals using~\eqref{eq:ci}
\State validate increment variance/correlation, finiteness, and positivity
\end{algorithmic}
\end{algorithm}

\subsection{Single-file code structure}
The requested submission has two primary source files.  Accordingly, the code
is organised as small, model-agnostic functions inside
\texttt{financial\_sde\_project.py}, rather than split across modules.

\begin{table}[H]
\centering
\small
\begin{tabular}{p{0.34\textwidth}p{0.58\textwidth}}
\toprule
\textbf{Function group} & \textbf{Purpose}\\
\midrule
\texttt{gbm\_terminal\_coupled}, \texttt{run\_gbm}
& Exact/EM/Milstein coupling, strong and weak estimators, positivity and matched-cost results.\\
\texttt{logistic\_vol}, \texttt{em\_nonaffine\_from\_increments}
& Stable model coefficients and swappable EM integration from supplied increments.\\
\texttt{aggregate\_increments}, \texttt{run\_nonaffine}
& Nested Brownian coupling, reference comparison, call payoff and diagnostics.\\
\texttt{nonaffine\_reference\_refinement}
& Coupled $2048\rightarrow4096\rightarrow8192$ reference refinement and slope sensitivity.\\
\texttt{plot\_*}, \texttt{validate}
& Reproduce all six figures and execute ten PASS/FAIL checks.\\
\bottomrule
\end{tabular}
\caption{Code organisation.  Integrators receive increments explicitly, so the
random driver and numerical method can be tested independently.}
\label{tab:code-structure}
\end{table}

\subsection{Reproducible protocol}
The default run uses seed 20260918, $M=300{,}000$ GBM paths, $M=20{,}000$
non-affine paths, batches of at most 1000 paths, and $N_{\rm ref}=2048$.  The
command
\begin{center}
\texttt{python financial\_sde\_project.py}
\end{center}
regenerates \texttt{generated/results.json},
\texttt{run\_output.txt}, and Figures~\ref{fig:gbm-convergence}--\ref{fig:reference-sensitivity}
with fixed names and at 220 DPI.  No absolute paths are used.  A reduced
\texttt{--quick} mode is available only as a smoke test; all report numbers are
from the default run.  NumPy performs vectorised path operations
\cite{harris2020}, and Matplotlib produces the figures \cite{hunter2007}.

\subsection{Oracle chain and automatic checks}
Two independent GBM oracles are used: the pathwise exact transition and the
closed-form terminal moments/distribution.  The non-affine reference is compared
with two finer grids on a fresh, commonly coupled validation sample.  Ten automatic checks cover expected-order
bands, the exact GBM mean, both increment variances, target correlation,
finiteness, positivity, and a plausible non-affine strong slope.  All ten pass
in the archived output.  The checks are deliberately broad enough to avoid
passing only one random realisation, yet narrow enough to catch a wrong
increment scaling, correlation formula, or scheme.

\subsection{Declared cost measure}
The cost measure is the number of time steps per path, equivalently one
drift/diffusion coefficient evaluation per step in the implemented scalar GBM
loop.  The comparison excludes plotting, random-number setup, file I/O,
compilation, and vector-allocation overhead.  It also does not pretend that an
EM and a Milstein step have identical floating-point operation counts: Milstein
evaluates an extra correction.  The reported ratio is therefore an algorithmic
step-count comparison, not a wall-clock speedup claim.

% ============================================================
\section{Numerical experiments and results}\label{sec:results}
% ============================================================

\subsection{GBM strong and weak convergence}
Table~\ref{tab:gbm-convergence} reports the mesh range used for the
fit.  Figure~\ref{fig:gbm-convergence} overlays the predicted slopes and 95\%
fit bands.  The observed EM slope $0.4899\pm0.0026$ is close to $1/2$; the
Milstein slope $0.9099\pm0.0167$ approaches order one.  The small shortfall from
one is consistent with a finite pre-asymptotic mesh range rather than evidence
for a different theoretical order.  The analytic mean-bias slope
$0.9982\pm0.0004$ is essentially one.

\begin{table}[H]
\centering
\small
\begin{tabular}{rrrrr}
\toprule
$N$ & $h$ & EM strong MAE & Milstein strong MAE & exact mean bias magnitude\\
\midrule
4   & 0.250000 & 4.52887 & 0.460308 & 0.0325759\\
8   & 0.125000 & 3.28870 & 0.268833 & 0.0163567\\
16  & 0.062500 & 2.34715 & 0.149881 & 0.00819567\\
32  & 0.031250 & 1.67424 & 0.0800666 & 0.00410218\\
64  & 0.015625 & 1.18565 & 0.0414826 & 0.00205218\\
128 & 0.0078125 & 0.838329 & 0.0211342 & 0.00102636\\
256 & 0.00390625 & 0.594648 & 0.0106906 & 0.000513248\\
\bottomrule
\end{tabular}
\caption{GBM convergence with $300{,}000$ coupled paths.  Error decreases at
the theoretically expected rates; the exact weak-bias column is common to EM
and Milstein for the linear payoff $\varphi(s)=s$.}
\label{tab:gbm-convergence}
\end{table}

\begin{figure}[H]
\centering
\includegraphics[width=0.98\textwidth]{fig01_gbm_convergence.png}
\caption{Coupled GBM convergence.  Left: EM and Milstein terminal mean absolute
errors follow the $h^{1/2}$ and $h$ reference trends, respectively; fit bands
make the slope uncertainty visible.  Right: the analytically evaluable mean
bias follows order one.  This separates discretisation error from Monte Carlo
noise and directly supports the claimed orders.}
\label{fig:gbm-convergence}
\end{figure}

\subsection{Matched-accuracy cost comparison}
At the predeclared target
$\mathbb E|S_N-S_T|\le0.85$, the coarsest tested EM grid that passes uses
$N=128$ and attains 0.83833; Milstein passes at $N=2$ with 0.75433.
Under the stated step-count measure, the work ratio is therefore
\begin{equation}
 \frac{C_{\rm EM}}{C_{\rm Mil}}=\frac{128}{2}=64.
\end{equation}
This is a valid matched-accuracy comparison, but it is not a claim of a
64-fold wall-clock speedup.  The negative result implicit in the comparison is
also important: at very loose accuracy, EM's cheaper algebra could matter;
that regime is not measured by this target.

\begin{figure}[htbp]
    \centering
    \includegraphics[width=0.82\textwidth]
    {generated/fig07_gbm_cost_accuracy.png}
    \caption{GBM strong error versus the declared computational cost, measured as
    the number of time steps per path. Both methods are compared at the same target
    mean absolute terminal error of \(0.85\) price units. Euler--Maruyama first meets
    the target at \(N=128\), whereas Milstein first meets it at \(N=2\), giving a
    \(64{:}1\) step-count ratio. This is an algorithmic step-count comparison, not a
    claim of a \(64\)-fold wall-clock speedup; random-number generation, plotting,
    file I/O, setup, and the extra arithmetic in a Milstein step are not included.}
    \label{fig:gbm-cost-accuracy}
\end{figure}

\subsection{Terminal distribution and moment validation}
The exact formulas give
\begin{equation}
 \mathbb E[S_T]=105.127110,\qquad
 \operatorname{Var}(S_T)=1917.591686.
\end{equation}
The exact-transition Monte Carlo sample gives 105.088894 and 1915.728882,
respectively; the mean is within four Monte Carlo standard errors, satisfying
the automatic oracle check.  At $N=64$, the 1\%, 5\%, 50\%, 95\%, and 99\%
quantiles are shown in Table~\ref{tab:quantiles}.

\begin{table}[H]
\centering
\small
\begin{tabular}{lrrrrr}
\toprule
Method & 1\% & 5\% & 50\% & 95\% & 99\%\\
\midrule
Exact lognormal & 38.2687 & 50.2610 & 97.0446 & 187.3749 & 246.0930\\
EM, $N=64$ & 37.9146 & 50.0592 & 97.0464 & 187.5114 & 244.2812\\
Milstein, $N=64$ & 38.3043 & 50.2674 & 96.9520 & 187.6476 & 245.1338\\
\bottomrule
\end{tabular}
\caption{GBM terminal quantiles in price units.  Central quantiles agree
closely; the tails reveal the remaining finite-step error more sensitively.}
\label{tab:quantiles}
\end{table}

\begin{figure}[H]
\centering
\includegraphics[width=0.98\textwidth]{fig02_gbm_distribution.png}
\caption{Terminal-law validation at $N=64$.  The density overlay checks the
whole distribution rather than only two moments, while the quantile plot makes
tail discrepancies visible.  Both methods approximate the exact lognormal
quantiles; Milstein is closer at some, but not all, reported quantiles.}
\label{fig:gbm-distribution}
\end{figure}

\subsection{Positivity stress test}
At $h=1$, 0.435\% of EM paths encounter a non-positive step multiplier; at
$h=1/2$, the rate is 0.034\%, and no violation is observed among 300,000 paths
for $h\le1/4$.  This is consistent with the rapid Gaussian-tail decay in
\eqref{eq:em-neg-prob}.  No Milstein violation occurs for the tested parameters,
as predicted by~\eqref{eq:milstein-min}.

\begin{figure}[H]
\centering
\includegraphics[width=0.72\textwidth]{fig03_gbm_positivity.png}
\caption{Observed fraction of paths with at least one non-positive update
multiplier.  EM fails on coarse grids even though the true GBM remains positive;
zeros at finer grids mean ``none in 300,000 paths'', not a mathematical
guarantee.  Milstein is positive here because its quadratic multiplier has a
positive minimum for the tested $h$, not because Milstein is unconditionally
positivity preserving.}
\label{fig:gbm-positivity}
\end{figure}

\subsection{Non-affine nested-grid convergence}
Table~\ref{tab:nonaffine} compares each coarse grid with the coupled 2048-step
reference.  Fitting the first four rows (the final row is reserved as a
near-reference diagnostic) gives strong slope $0.5111\pm0.0047$, closely
matching the EM benchmark $1/2$.  Every 95\% confidence interval for the paired
call-payoff difference contains zero.  The correct conclusion is therefore
that sampling uncertainty is already comparable to the estimated weak bias;
the data do not support quoting a weak slope.

\begin{table}[H]
\centering
\small
\begin{tabular}{rrrrr}
\toprule
$N$ & $h$ & strong MAE & strong 95\% CI & call bias (95\% CI)\\
\midrule
16  & 0.062500 & 0.886321 & [0.875404, 0.897238] & 0.004204 $[-0.010095,0.018504]$\\
32  & 0.031250 & 0.627466 & [0.619890, 0.635042] & 0.003974 $[-0.006056,0.014005]$\\
64  & 0.015625 & 0.441095 & [0.435856, 0.446335] & 0.003689 $[-0.003286,0.010664]$\\
128 & 0.0078125 & 0.306001 & [0.302342, 0.309660] & 0.002141 $[-0.002714,0.006996]$\\
256 & 0.00390625 & 0.210677 & [0.208183, 0.213171] & 0.002748 $[-0.000568,0.006064]$\\
\bottomrule
\end{tabular}
\caption{Non-affine EM convergence against the coupled $N_{\rm ref}=2048$
numerical reference using 20,000 paths.  Strong convergence is resolved, while
the weak payoff bias is not statistically distinguishable from zero.}
\label{tab:nonaffine}
\end{table}

\begin{figure}[H]
\centering
\includegraphics[width=0.98\textwidth]{fig04_nonaffine_convergence.png}
\caption{Nested-grid evidence for the non-affine model.  Left: the pathwise
terminal error follows an empirical half-order trend with narrow Monte Carlo
intervals.  Right: all payoff-bias confidence intervals meet zero, so claiming
a weak order from these samples would over-interpret noise.}
\label{fig:nonaff-convergence}
\end{figure}

\subsection{Increment, finiteness, positivity, and reference checks}
The fine-grid increment diagnostics are
\begin{equation}
 \frac{\widehat{\operatorname{Var}}(\Delta W^{(1)})}{h}=1.00526,
 \qquad
 \frac{\widehat{\operatorname{Var}}(\Delta W^{(2)})}{h}=1.00346,
 \qquad
 \widehat{\operatorname{Corr}}=-0.70185,
\end{equation}
against targets $(1,1,-0.7)$.  All 20,000 reference paths are finite and
positive.  Their estimated mean terminal price is $105.01815$ with standard
error $0.21433$; the undiscounted call-payoff mean estimate is $14.48371$ with standard error
$0.15090$.  These are model outputs, not closed-form benchmarks.

On 4000 fresh validation paths (seed 20269918), Brownian increments are
generated at 8192 steps and aggregated to all coarser grids.  The 2048-,
4096-, and 8192-step references therefore share the same Brownian paths.
Refining from 2048 to 4096 steps gives a mean absolute terminal difference of
\begin{equation}
 0.057485\quad\text{with 95\% CI }[0.055953,0.059016].
 \label{eq:refinement-result}
\end{equation}
This value is recomputed on the common three-grid sample and supersedes the
earlier two-grid validation value; the main 20,000-path experiment is unchanged.

\begin{table}[H]
\centering
\begin{tabular}{rrr}
\toprule
Reference steps & Terminal MAE & 95\% MC CI\\
\midrule
$2048\rightarrow4096$ & 0.057485 & [0.055953, 0.059016]\\
$4096\rightarrow8192$ & 0.040731 & [0.039653, 0.041809]\\
\bottomrule
\end{tabular}
\caption{Two reference refinements on the same 4000 validation paths.
These are differences between numerical solutions, not exact-solution errors.}
\label{tab:reference-refinement}
\end{table}

The second mean difference is 0.708551 times the first, compared with
$1/\sqrt{2}\approx0.7071$ as a half-order scaling benchmark.
The decreasing differences support refinement consistency for this experiment;
they do not provide an upper bound for the error relative to the exact solution.

\paragraph{Sensitivity of the fitted slope to the reference grid.}
For each reference, the same 4000 validation paths are also integrated at
$N=16,32,64,128,256$.  All three fits use $N=16,32,64,128$, retaining
$N=256$ only as a diagnostic, as in the main experiment.
\begin{table}[H]
\centering
\begin{tabular}{rrr}
\toprule
$N_{\rm ref}$ & Fitted slope & OLS standard error\\
\midrule
2048 & 0.501554 & 0.006405\\
4096 & 0.492814 & 0.006792\\
8192 & 0.489893 & 0.006917\\
\bottomrule
\end{tabular}
\caption{Reference-grid sensitivity with fixed validation paths and fitting
range.  Standard errors describe the log--log fits, not confidence intervals
accounting for all Monte Carlo and reference uncertainty.}
\label{tab:reference-slopes}
\end{table}
Relative to the 2048-step validation fit, the maximum slope change is 0.011662.
The three slopes remain close to one half, supporting the empirical half-order
trend under reference refinement over this fitting range.  These validation
slopes are separate from the main 20,000-path estimate $0.5111$; differences
between samples should not be attributed solely to the reference grid.
The full gridwise MAEs and confidence intervals are recorded in
\texttt{generated/results.json}.  This study checks two finer reference step
sizes, but the 8192-step solution remains numerical and weak payoff bias is
still unresolved.

\begin{figure}[H]
\centering
\includegraphics[width=0.98\textwidth]{fig05_nonaffine_paths.png}
\caption{Representative non-affine paths.  Integrating $X=\log S$ guarantees
positive displayed prices after exponentiation (left), while the logistic
map keeps instantaneous volatility strictly between 0.10 and 0.50
year$^{-1/2}$ (right).  These structural properties are enforced by the model
representation rather than inferred from a finite sample.}
\label{fig:nonaff-paths}
\end{figure}

\begin{figure}[H]
\centering
\includegraphics[width=0.98\textwidth]{fig06_reference_sensitivity.png}
\caption{Supplementary reference validation on 4000 common paths.
Left: adjacent-reference differences with 95\% Monte Carlo intervals.
Right: coarse-grid differences against three references; the legend slopes
use only $N=16,32,64,128$.  The empirical trend persists under refinement.}
\label{fig:reference-sensitivity}
\end{figure}

% ============================================================
\section{Limitations and robustness}\label{sec:limitations}
% ============================================================
\begin{enumerate}[leftmargin=2.2em]
  \item \textbf{Numerical reference.}  The non-affine error is relative to EM
  with 2048 steps.  Equation~\eqref{eq:refinement-result} measures a difference
  between two numerical references, not the error relative to the exact
  solution.  Two finer references and the fixed-sample slope comparison
  support the observed trend, but do not eliminate reference error.
  \item \textbf{Weak convergence resolution.}  Paired common random numbers
  reduce variance, yet the call-bias intervals still contain zero.  The report
  therefore does not manufacture a weak slope.  More paths, multilevel Monte
  Carlo, or a different payoff would be needed.
  \item \textbf{Regression uncertainty.}  Slope standard errors diagnose the
  chosen log--log fit.  They are not theorem proofs and do not include every
  model-selection uncertainty, such as mesh-range choice.
  \item \textbf{Cost definition.}  Step count is portable but omits the extra
  arithmetic in Milstein and hardware effects.  A wall-clock study should use
  repeated timings on the same machine and separately count random draws and
  coefficient/derivative evaluations.
  \item \textbf{Model risk.}  Numerical convergence does not validate GBM or
  the non-affine model as a description of market data.  Calibration,
  risk-neutral drift, transaction costs, jumps, and empirical fit are outside
  this numerical-analysis project.
  \item \textbf{Extreme paths.}  The stable logistic evaluation and batched
  finiteness checks reduce numerical risk, but $e^X$ can overflow under much
  longer horizons or more extreme parameters.  Production code should retain
  log values and define an overflow policy.
\end{enumerate}

% ============================================================
\section{Conclusions}\label{sec:conclusion}
% ============================================================
\begin{itemize}[leftmargin=2em]
  \item Exact GBM coupling recovered the expected strong orders: EM
  $0.4899\pm0.0026$ and Milstein $0.9099\pm0.0167$; the mean weak bias had order
  $0.9982\pm0.0004$.
  \item At matched strong MAE $\le0.85$, Milstein required 2 steps versus 128
  for EM under the declared step-count cost, while the report avoids converting
  that 64:1 ratio into an unsupported wall-clock claim.
  \item Distribution and moment checks agree with the exact lognormal oracle,
  and the coarse-grid stress test shows why price-space EM is not
  positivity-preserving.
  \item For the non-affine model, nested Brownian coupling produced empirical
  half-order convergence relative to the 2048-step reference and passed the
  reported variance, correlation, terminal-finiteness, and terminal-positivity
  checks.  Two reference refinements and a common-sample slope comparison
  support stability of the half-order trend over the tested fitting range.
  Weak payoff bias remained
  unresolved at the available Monte Carlo precision.
  \item A natural extension is multilevel Monte Carlo, which would reuse the
  same nested coupling to reduce the cost of resolving weak payoffs
  \cite{giles2008}; a second extension is a correctly implemented
  multidimensional Milstein scheme with iterated stochastic integrals.
\end{itemize}

% ============================================================
\addcontentsline{toc}{section}{References}
% ============================================================
\begin{thebibliography}{99}
\bibitem{higham2001}
D. J. Higham, ``An algorithmic introduction to numerical simulation of
stochastic differential equations,'' \emph{SIAM Review}, vol.~43, no.~3,
pp.~525--546, 2001. doi: \href{https://doi.org/10.1137/S0036144500378302}{10.1137/S0036144500378302}.

\bibitem{kloedenplaten1992}
P. E. Kloeden and E. Platen, \emph{Numerical Solution of Stochastic
Differential Equations}. Berlin: Springer, 1992.
doi: \href{https://doi.org/10.1007/978-3-662-12616-5}{10.1007/978-3-662-12616-5}.

\bibitem{milstein1995}
G. N. Milstein, \emph{Numerical Integration of Stochastic Differential
Equations}. Dordrecht: Springer, 1995.
doi: \href{https://doi.org/10.1007/978-94-015-8455-5}{10.1007/978-94-015-8455-5}.

\bibitem{giles2008}
M. B. Giles, ``Multilevel Monte Carlo path simulation,''
\emph{Operations Research}, vol.~56, no.~3, pp.~607--617, 2008.
doi: \href{https://doi.org/10.1287/opre.1070.0496}{10.1287/opre.1070.0496}.

\bibitem{harris2020}
C. R. Harris et al., ``Array programming with NumPy,'' \emph{Nature},
vol.~585, pp.~357--362, 2020.
doi: \href{https://doi.org/10.1038/s41586-020-2649-2}{10.1038/s41586-020-2649-2}.

\bibitem{hunter2007}
J. D. Hunter, ``Matplotlib: A 2D graphics environment,''
\emph{Computing in Science \& Engineering}, vol.~9, no.~3, pp.~90--95, 2007.
doi: \href{https://doi.org/10.1109/MCSE.2007.55}{10.1109/MCSE.2007.55}.

\bibitem{coursebrief}
MTH321 teaching team, \emph{Project 1 Brief: Numerical Solution of ODE Initial
Value Problems}, extension pathway 4, 2026, course material supplied with the
project.

\bibitem{problempack}
MTH321 teaching team, \emph{Problem Pack: Five Suggested Directions for the ODE
IVP Project}, ``Financial SDE: GBM benchmark and non-affine volatility,'' 2026,
course material supplied with the project.
\end{thebibliography}

\appendix

% ============================================================
\section{Reproduction record and complete console output}
% ============================================================
\subsection{Environment and commands}
Install the declared dependencies in any isolated Python environment:
\begin{center}
\texttt{python -m pip install numpy matplotlib}
\end{center}
From the repository root run
\begin{center}
\texttt{python financial\_sde\_project.py}.
\end{center}
All paths are relative; the output directory is created automatically.  The
JSON file records Python, NumPy, Matplotlib, platform, seed, sample sizes,
parameters, every table row, validation results, and wall time.

\subsection{Archived program output}
\lstinputlisting[language={},caption={Complete output from the full reproducibility run.}]{generated/run_output.txt}

% ============================================================
\section{Complete Python source code}
% ============================================================
The following listing is the exact code used to generate every reported result
and figure.
\IfFileExists{../financial_sde_project.py}{%
\lstinputlisting[caption={Complete source: \texttt{financial\_sde\_project.py}.}]{../financial_sde_project.py}%
}{%
\lstinputlisting[caption={Complete source: \texttt{financial\_sde\_project.py}.}]{financial_sde_project.py}%
}

% ============================================================
\section{AI Transparency Log}
% ============================================================
\subsection{Tools used}
\begin{table}[H]
\centering
\small
\begin{tabular}{p{0.25\textwidth}p{0.20\textwidth}p{0.46\textwidth}}
\toprule
Tool & Version/date & Primary use\\
\midrule
OpenAI Codex & 18 September 2026 & Requirements mapping, derivation drafting,
Python implementation, experiment execution, result interpretation, and
LaTeX assembly.\\
NumPy / Matplotlib & Versions recorded in \texttt{results.json} & Vectorised
Monte Carlo computation and deterministic figure generation.\\
\bottomrule
\end{tabular}
\caption{AI and software tools.  Team members must replace or supplement this
record if they make later changes with other tools.}
\end{table}

\subsection{Significant AI-assisted claim audits}

% C.2 中替换表格的开始和结束
\begingroup
\scriptsize
\setlength{\tabcolsep}{4pt}
\renewcommand{\arraystretch}{1.15}

\begin{tabular}{|p{0.05\textwidth}|p{0.44\textwidth}|p{0.44\textwidth}|}
\hline
\textbf{Item} &
\textbf{Prompt (near-verbatim) and assumptions} &
\textbf{Risk and possible falsifier} \\
\hline

% 保留上一版给你的第 1 行内容
1 &
``Please derive the Euler--Maruyama and Milstein updates for GBM and help me build a NumPy convergence test. Use the same Brownian increments for EM, Milstein, and the exact GBM path, then calculate the terminal mean absolute strong error for several values of \(N\) and fit the log--log slopes. I expect orders \(1/2\) and \(1\), but please include checks that could reveal an implementation mistake.'' &
The slope could look convincing even if the simulation is wrong. In particular, using \(Z\) instead of \(\sqrt{h}Z\), generating a separate exact path, or fitting errors before the asymptotic range would change the result. We would reject the claim if the exact-path mean disagreed with \(S_0e^{\mu T}\), if errors did not decrease consistently, or if the fitted slopes fell outside the automatic acceptance bands. \\
\hline

% 保留上一版给你的第 2 行内容
2 &
``For the two-factor non-affine model there is no exact solution available. Can we use a 2048-step EM solution as a numerical reference and compare the coarser grids with it? Please construct every grid from the same fine Brownian increments, check the result again with 4096 and 8192 steps, and make sure the report does not describe the reference as exact. Also, do not claim a weak convergence rate if the payoff-bias confidence intervals still cross zero.'' &
A 2048-step path still contains discretisation error, so measuring against it can make the coarse-grid error appear smaller and can distort the fitted slope. There is also a coding risk that the coarse increments are regenerated instead of summed from the fine increments. The claim would be weakened if the \(2048\mathbin{\to}4096\) and \(4096\mathbin{\to}8192\) differences failed to decrease, if the fitted slope changed substantially when the reference was refined, or if the measured increment variance and correlation missed their targets \(1\), \(1\), and \(-0.7\). \\
\hline

% 保留上一版给你的第 3 行内容
3 &
``Using the GBM results, compare EM and Milstein at the same strong-error target rather than at the same number of steps. Use a target MAE of \(0.85\), find the coarsest tested grid that passes for each method, and report the ratio. Is it valid to call that ratio a speedup? Please state exactly what is being counted and what is excluded.'' &
The comparison counts time steps, but one Milstein step includes the additional quadratic correction and therefore costs more arithmetic than one EM step. Random-number generation, array allocation, and hardware effects are also absent from the count. The \(64{:}1\) result would not justify a \(64\)-fold runtime claim; repeated timings or an operation count at the same MAE could produce a smaller advantage or, in another accuracy regime, reverse the ranking. \\
\hline
\end{tabular}

\endgroup

\subsection{Independent verification and decisions}

\begingroup
\scriptsize
\setlength{\tabcolsep}{4pt}
\renewcommand{\arraystretch}{1.15}

\begin{tabular}{|p{0.05\textwidth}|p{0.44\textwidth}|p{0.44\textwidth}|}
\hline
\textbf{Item} &
\textbf{Independent check and evidence location} &
\textbf{Final decision, modification, or limitation} \\
\hline

1 &
We checked the simulated paths against the exact GBM transition in Equation~(13), the exact mean and variance in Equation~(4), and the analytic mean-bias formula in Equation~(18). The errors in Table~2 decrease across the full mesh sequence, and the automatic tests check the fitted-order bands and the exact terminal mean. &
We retained the strong-order conclusions because the fitted slopes, \(0.4899\) for EM and \(0.9099\) for Milstein, agree with the expected behaviour over the tested grids. The weak order is supported by the analytic bias slope \(0.9982\), not inferred only from a noisy Monte Carlo mean. The reported slope standard errors are described as regression diagnostics rather than proofs of the theoretical orders. \\

\addlinespace

2 &
A separate validation run used 4000 fresh paths with seed \(20269918\). Increments were generated on the 8192-step grid and summed to construct the 4096-, 2048-, and coarser grids. The adjacent-reference MAEs decrease from \(0.057485\) to \(0.040731\) in Table~5, while the fitted slopes against the three references remain between \(0.489893\) and \(0.501554\) in Table~6. Increment variance, correlation, finiteness, and terminal positivity were also checked. &
We kept the half-order statement but called it an \emph{empirical} trend relative to a numerical reference. We did not call the 2048-step solution exact, and we retained reference error as a limitation. Because every call-payoff bias interval in Table~4 contains zero, we did not calculate or report a weak payoff slope. \\

\addlinespace

3 &
From the stored pathwise errors, we reapplied the same requirement,
\(\mathbb{E}|S_N-S_T|\leq 0.85\), to both schemes. The first tested grids meeting it were \(N=128\) for EM, with MAE \(0.83833\), and \(N=2\) for Milstein, with MAE \(0.75433\). Section~4.5 separately lists the operations excluded from the cost measure. &
We reported \(128/2=64\) only as a time-step-count ratio at the stated error target. We removed any interpretation of this value as a measured runtime speedup. Establishing an actual speedup would require repeated wall-clock timings on the same machine or a more detailed count of arithmetic, coefficient evaluations, and random-number costs. \\

\hline
\end{tabular}

\endgroup

\subsection{Intellectual ownership declaration}
Before submission, every member must complete and sign the table below.  Each
signer affirms that they can derive the equations, explain and modify the
team-owned code, identify trusted library interfaces, and reproduce the
validation without AI assistance during the viva.

\begin{table}[H]
\centering
\small
\begin{tabular}{p{0.17\textwidth}p{0.15\textwidth}p{0.35\textwidth}p{0.24\textwidth}}
\toprule
Name & Signature & Team-owned artefact/contribution & Trusted dependency/interface\\
\midrule
\texttt{Jinyang Wang} & & & \\
\texttt{Yangxiaoxiang Zheng} & & & \\
\texttt{Jinsen Chen} & & & \\
\texttt{Zheyuan Zheng} & & & \\
\texttt{Dongdong Hao} & & & \\
\bottomrule
\end{tabular}
\end{table}

% ============================================================
\section{Individual Contribution Statements}
% ============================================================

\begin{formbox}
\textbf{Individual Contribution Statement --- Project 1 (Financial SDE)}\\[0.4em]

\textbf{Name:} \texttt{Jinyang Wang}
\hfill
\textbf{Role:} \texttt{Project Manager}\\[0.4em]

\textbf{The specific things I did this project:}

\begin{enumerate}[leftmargin=2em]

\item I established the team's working structure and coordinated its weekly
progress. In Week 1, I created the Git repository with the required
\path{code/}, \path{figures/}, and \path{report/} structure, added all five
team members and the instructor as collaborators, and confirmed that every
member had cloned the repository and pushed at least once. I then ran the
weekly team meetings and maintained the Gantt and milestone records for
Weeks 1--3, assigning named tasks and tracking blockers. These included the
Week 1 issue of evaluating the logistic-volatility function stably and the
Week 2 issue of distinguishing a numerical non-affine reference from an
exact solution.

\item I coordinated the mathematical direction and reporting strategy. I led
the selection of Topic ④ and prepared the defence used in Sections 1.1--1.2:
GBM was selected because its exact transition provides a pathwise oracle for
separating discretisation error from Monte Carlo uncertainty, while the
non-affine extension requires the team to construct and validate a numerical
reference. I used this distinction to organise responsibilities across the
mathematical-theory, algorithm, validation, and visualization roles and to
ensure that each experiment had a clear evidential purpose.

\item I coordinated final integration and independently checked the
reproducibility claim in Section 4.3. I collected the report sections from
their respective owners, ran
\path{python financial_sde_project.py}, and compared the reproduced console
output and \path{generated/results.json} with the values reported in
Tables 2--6 and Appendix A.2. I also compiled the team AI Transparency Log in
Appendix C from the members' audit notes and coordinated completion of the
five signed ICS forms in Appendix D.

\end{enumerate}

\textbf{The hardest conceptual obstacle and how I overcame it:}\\
The hardest obstacle was converting a broad project brief into parallel,
traceable work while keeping the distinction between exact and numerical
validation clear. I overcame this by assigning an owner to every deliverable,
recording blockers in the weekly milestones, and using GBM's exact oracle and
the non-affine nested numerical reference as two separate validation chains.
This allowed the team to work in parallel without treating the numerical
reference as an exact answer.\\[0.4em]

\textbf{One AI-assisted item I used and what I changed:}\\
I cross-reference AI Transparency Log \textbf{Item 2}, concerning the nested
Euler--Maruyama reference for the non-affine model. AI assistance helped us
identify the reference-construction and refinement requirements, but I
required the team to change the wording from ``exact reference'' to
``numerical reference,'' add an independent refinement check, and state the
measured strong rate as empirical. I also ensured that the associated risk,
falsifier, evidence location, and final limitation were recorded in
Appendix C rather than accepting the AI-generated claim without an audit.\\[0.4em]

\textbf{What I would do differently / what I still need to understand:}\\
I would study the project brief, assessment rubric, and ICS requirements in
detail before the first meeting and prepare the repository structure, task
allocation, submission checklist, and weekly schedule immediately. Because I
adjusted my classes during the first week, I did not organise the project as
early as I should have, which caused later tasks to be completed under time
pressure. I also should have asked for help interpreting the ICS criteria
before drafting my own version, rather than revising it after discovering that
the contribution claims needed to be more specific and traceable.\\[0.6em]

\textbf{Signature:} \underline{\hspace{4cm}}\qquad
\textbf{Date:} \underline{\hspace{2.5cm}}

\end{formbox}

\newpage



\begin{formbox}
\textbf{Individual Contribution Statement --- Project 1 (Financial SDE)}\\[0.4em]

\textbf{Name:} \texttt{Yangxiaoxiang Zheng (1929725)}
\hfill
\textbf{Role:} \texttt{Mathematical Theory}\\[0.4em]

\textbf{The specific things I did this project:}

\begin{enumerate}[leftmargin=2em]

\item I completed and checked the mathematical derivations presented in
Sections 2--3 and Equations (1)--(29). These included the It\^o
log-transformation and exact solution of GBM, its terminal moments, the
non-affine stochastic-volatility model, the Euler--Maruyama and Milstein
schemes, pathwise coupling, strong and weak error definitions, the weak-bias
expansion, positivity conditions, correlated Brownian increments, nested-grid
coupling, and Monte Carlo uncertainty quantification. I also checked that the
assumptions and limitations stated with these formulas were consistent with
the Financial SDE extension pathway.

\item I derived the Milstein correction in Equations (11)--(12). Setting
$Y_s=W_s-W_{t_n}$ and applying It\^o's formula gives
\[
  d(Y_s^2)=2Y_s\,dW_s+ds,
\]
and hence
\[
  \int_{t_n}^{t_{n+1}}\int_{t_n}^{s}dW_r\,dW_s
  =\frac{1}{2}\left[(\Delta W_n)^2-h\right].
\]
Substituting $b(S)=\sigma S$ and $b'(S)=\sigma$ produces the scalar GBM
Milstein update. I checked the factor $1/2$, the $-h$ term, and the distinction
between It\^o and Stratonovich integration. I also established the limitation
stated in Section 3.6: the scalar Milstein formula cannot be transferred
directly to the two-noise non-affine model because its diffusion fields
produce cross derivatives and multidimensional iterated stochastic integrals.

\item I independently verified the derivations and checked their consistency
with the implementation. I re-derived the Milstein double integral using
It\^o's formula, checked the strong/weak error definitions and weak-bias
expansion through conditional expectations, and verified the EM and Milstein
positivity conditions. I then compared the mathematical update formulas with
the algorithm implementation to ensure that the shared Brownian increments,
Milstein correction, correlated increments, and nested-grid sums matched the
report. Evidence is located in Sections 3.2--3.8 and AI Transparency Log
Items 1--3.

\end{enumerate}

\textbf{The hardest conceptual obstacle and how I overcame it:}\\
The hardest conceptual obstacle was understanding why the Milstein correction
contains both the factor $1/2$ and the term $-h$, and why the resulting scalar
formula is valid for GBM but not automatically for the non-affine two-noise
system. I overcame this by deriving the iterated stochastic integral directly
from It\^o's formula, checking every substitution into the GBM diffusion
coefficient, and then examining the cross derivatives of the non-affine
diffusion fields before deciding that a scalar Milstein extension would be
mathematically unjustified.\\[0.4em]

\textbf{One AI-assisted item I used and what I changed:}\\
I cross-reference AI Transparency Log \textbf{Item 1}, concerning the strong
orders of Euler--Maruyama and Milstein under pathwise coupling. AI assistance
was used to explain the It\^o--Taylor derivation and the double stochastic
integral. I did not accept the derivation unchanged: I independently
re-derived the integral using It\^o's formula, checked the $1/2$ and $-h$
terms, verified the conditional-mean calculations, and added the explicit
limitation that scalar Milstein cannot be applied directly to the non-affine
two-noise model.\\[0.4em]

\textbf{What I would do differently / what I still need to understand:}\\
I would begin earlier with the basic It\^o-calculus identities and organise
the theory in dependency order: exact GBM, Euler--Maruyama, scalar Milstein,
error definitions, positivity, and finally the multidimensional non-affine
model. I would also maintain a derivation-to-code checklist while the
implementation was being developed, rather than performing most consistency
checks near the end. I can independently explain the formulas used in this
project; a natural topic for further study is the construction of full
multidimensional Milstein schemes and their cross iterated integrals.\\[0.6em]

\textbf{Signature:} \underline{\hspace{4cm}}\qquad
\textbf{Date:} \underline{\hspace{2.5cm}}

\end{formbox}

\newpage




\begin{formbox}
\textbf{Individual Contribution Statement --- Project 1 (Financial SDE)}\\[0.4em]

\textbf{Name:} \texttt{Jinsen Chen (2363993)}
\hfill
\textbf{Role:} \texttt{Algorithm Implementation}\\[0.4em]

\textbf{The specific things I did this project:}

\begin{enumerate}[leftmargin=2em]

\item I developed the stochastic-algorithm layer recorded in two traceable
Git histories. On branch \texttt{jooshuua-algorithm}, commit
\texttt{e9e880f} added \texttt{Personal\_Algorithm\_Work.zip}, containing
\texttt{gbm.py}, \texttt{nonaffine.py}, \texttt{convergence.py}, and
\texttt{parameters.py}. On \texttt{main}, commit \texttt{dc1674a} added the
618-line \texttt{financial\_sde\_algorithm\_flow.py}. My functions include
\texttt{gbm\_exact\_step()}, \texttt{gbm\_em\_step()},
\texttt{gbm\_milstein\_step()}, \texttt{coupled\_gbm\_terminal()},
\texttt{correlated\_brownian\_increments()},
\texttt{nonaffine\_em\_terminal()}, and \texttt{aggregate\_increments()}.
They implement shared-increment GBM comparisons, correlated two-factor
noise, log-price Euler--Maruyama, and nested coarse/fine Brownian paths.

\item I implemented and reviewed algorithm-level validation. Evidence is in
\texttt{check\_algorithms.py} inside the zip from commit \texttt{e9e880f}.
It checks zero-coefficient GBM, preservation of fine-to-coarse increment
sums, the target Brownian correlation, bounded logistic volatility, finite
positive non-affine prices, and the observed EM and Milstein strong orders.
In commit \texttt{dc1674a}, the functions \texttt{mean\_ci95()},
\texttt{fitted\_slope()}, and \texttt{refine\_nonaffine\_reference()} add
confidence intervals, log--log order estimates, and coupled
2048/4096/8192-step reference-sensitivity checks.

\item For the experiments reported in the \emph{Discretisation and
mathematical analysis} and \emph{Implementation and reproducibility}
sections, I decided that exact GBM, EM, and Milstein must share Brownian
increments, and that each non-affine coarse path must be constructed by
summing increments from its fine path. I did not apply the scalar Milstein
formula to the two-noise model because a valid multidimensional method
requires cross terms and iterated stochastic integrals. Its measured strong
rate is therefore reported as empirical rather than as an exact-reference
error.

\end{enumerate}

\textbf{The hardest conceptual obstacle and how I overcame it:}\\
The hardest issue was comparing stochastic discretisations on the same
Brownian path. Independently generated paths would mix Monte Carlo variation
with time-discretisation error. I overcame this by sharing every
$\Delta W$ in the GBM comparison and constructing each non-affine coarse
increment as a sum of increments from the common fine path.\\[0.4em]

\textbf{One AI-assisted item I used and what I changed:}\\
I cross-reference AI Transparency Log \textbf{Item 2}, concerning the nested
Euler--Maruyama reference. I did not treat the suggested fine-grid result as
exact. I explicitly coupled all grids, added 2048/4096/8192-step refinement
checks and confidence intervals, and described the observed strong rate as
empirical. I also retained Euler--Maruyama rather than accepting an
unjustified scalar Milstein extension to the two-noise system.\\[0.4em]

\textbf{What I would do differently / what I still need to understand:}\\
I would merge \texttt{jooshuua-algorithm} through a reviewed pull request
and use smaller commits for GBM, non-affine simulation, and convergence
analysis, rather than later uploading one complete file to \texttt{main}.
I still need a deeper understanding of multidimensional Milstein methods,
especially cross iterated integrals for non-commutative noise.\\[0.6em]

\textbf{Signature:} \underline{\hspace{4cm}}\qquad
\textbf{Date:} \underline{\hspace{2.5cm}}

\end{formbox}

\newpage

\begin{formbox}
\textbf{Individual Contribution Statement --- Project 1 (Financial SDE)}\\[0.4em]

\textbf{Name:} \texttt{Zheyuan Zheng (2362114)}
\hfill
\textbf{Role:} \texttt{Visualization \& Report}\\[0.4em]

\textbf{The specific things I did this project:}

\begin{enumerate}[leftmargin=2em]

\item I wrote and refined the figure-generation code in
\path{financial_sde_project.py}. The functions
\path{plot_gbm_convergence()} and \path{add_loglog_fit_band()} produce the
GBM convergence figure, including OLS slopes and standard errors obtained
with \path{np.polyfit(...,cov=True)}, theoretical $h^{1/2}$ and $h$ reference
slopes, filled markers for the fitted range, and pointwise 95\% OLS fit
bands. I also implemented the error-versus-cost figure
\path{fig07_gbm_cost_accuracy.png}, the non-affine convergence figure, and
the reference-sensitivity figure
\path{fig06_reference_sensitivity.png}. These figures show the matched-error
step-count comparison, empirical non-affine convergence, and the stability
of the fitted slope under reference-grid refinement.

\item I checked every figure caption against
\path{latex/generated/results.json} and \path{run_output.txt}. I verified
the fitted slopes and their reported uncertainties, the common target
$\mathrm{MAE}\leq0.85$, the $M=300{,}000$ paths behind the observation floor
$1/M=3.3\times10^{-6}$, and the adjacent-reference differences. The evidence
is located in the captions of Figures 1--7 and in those two generated files.
I also checked that the number of figures quoted in the report matched the
seven figures produced by the reproducible workflow.

\item I made three reporting decisions to prevent the figures from
overstating the evidence. In Section 5.4, a zero observed violation rate is
reported as ``none observed in 300,000 paths,'' accompanied by the
observation floor $1/M$ and the rule-of-three bound
$3/M=1.0\times10^{-5}$, rather than being presented as a mathematical
guarantee. In Section 5.5, the finest grid $N=256$ is excluded from the
slope fit as a near-reference diagnostic, and the fitted non-affine rate is
labelled empirical because the reference is numerical. In the
reference-sensitivity figure, the change from $0.5016$ to $0.4899$ is
reported as evidence that the approximate half-order trend is not an
artefact of one reference grid.

\end{enumerate}

\textbf{The hardest conceptual obstacle and how I overcame it:}\\
The hardest obstacle was distinguishing two different kinds of uncertainty.
A paired Monte Carlo confidence interval measures sampling noise, whereas an
OLS band around a log--log fit measures regression uncertainty conditional
on the selected grids. My first draft used one type of band and could have
suggested that slope uncertainty came entirely from randomness. I separated
the paired 95\% Monte Carlo intervals from the OLS fit bands and used the
analytic-bias panel as a noise-free control. A second difficulty was plotting
a zero observed rate on a logarithmic axis; I used the $1/M$ observation
floor and the $3/M$ rule-of-three bound to distinguish ``not observed'' from
``impossible.''\\[0.4em]

\textbf{One AI-assisted item I used and what I changed:}\\
I cross-reference AI Transparency Log \textbf{Item 3}, concerning the
matched-accuracy cost comparison. The AI-assisted draft plotted error
against step count and described Milstein as 64 times faster. I changed the
figure and interpretation by marking only the coarsest grids that passed the
same threshold, $\mathrm{MAE}\leq0.85$, and by defining cost explicitly as
time steps per path. I listed the excluded costs---setup, file I/O, plotting,
compilation, allocation, and the extra algebra within a Milstein step---and
replaced the wall-clock speed claim with a limited 64:1 step-count comparison.
I also stated the falsifier: repeated wall-clock or operation-count
measurements at the same error target could reverse the conclusion.\\[0.4em]

\textbf{What I would do differently / what I still need to understand:}\\
I would establish the visual conventions before producing the figures,
including the cost measure, the meaning of each uncertainty band, the fitted
grid range, and the wording used for a zero observed rate. I would also ask
earlier which error definition each figure was intended to display. I can
explain the scalar Milstein correction used for GBM, but I still need a
deeper understanding of its It\^o--Taylor derivation and of the cross terms
and iterated stochastic integrals required for a full multidimensional
Milstein method.\\[0.6em]

\textbf{Signature:} Zheyuan Zheng
\qquad
\textbf{Date:} 18 September 2026

\end{formbox}

\newpage

\begin{formbox}
\textbf{Individual Contribution Statement --- Project 1 (Financial SDE)}\\[0.4em]

\textbf{Name:} Dongdong Hao (2364312)
\hfill
\textbf{Role:} Testing \& Validation\\[0.4em]

\textbf{The specific things I did this project:}
\begin{enumerate}[leftmargin=2em]

\item I independently tested the submitted implementation by checking out the exact Git commit
\path{ef1880737f33764acc016832a8953057e1bc8108}
and running \path{Team_H_Code/code/run_all.py} locally using Python 3.13.7.
The program exited successfully with \texttt{OVERALL: PASS}, and I recorded the exit status and the results of all ten built-in validation checks. I also repeated the run and compared the console summaries to confirm that the reported results were reproducible.

\item I reviewed the reported GBM results against the acceptance criteria implemented in the code. This included the Euler--Maruyama and Milstein strong-convergence slopes of \(0.4899\) and \(0.9099\), the analytic weak-bias slope of \(0.9982\), the exact terminal-mean comparison, and the variance and target-correlation checks for the generated Brownian increments. I documented the evidence in \path{financial_sde_test_report.md}, with supporting output stored in \path{run_output.txt} and \path{results.json}.

\item I validated the non-affine experiment by checking that the reference terminal values were finite and positive, confirming the empirical strong slope of \(0.5111\), and reviewing the reference-refinement sequence \(2048 \rightarrow 4096 \rightarrow 8192\). The adjacent-reference MAE decreased from \(0.0575\) to \(0.0407\), supporting the stability of the numerical-reference construction. I also recorded that this was numerical evidence rather than proof that the finest reference was exact.

\end{enumerate}

\textbf{The hardest conceptual obstacle and how I overcame it:}\\
The main difficulty was distinguishing a successful numerical rerun from a complete reproducibility guarantee. A passing result can still depend on the checked-out commit, Python version, dependency versions, random seeds, and the precise acceptance criteria. I addressed this by testing an exact Git commit, recording the Python version and exit status, preserving the console output and JSON results, and checking each reported convergence or statistical value against the corresponding criterion. I also stated explicitly that my checks validated the numerical outputs, but did not establish the visual quality of every figure or reproduce every intermediate development state.\\[0.4em]

\textbf{One AI-assisted item I used and what I changed:}\\
For AI Transparency Log item 2, concerning the nested numerical reference for the non-affine model, I used Codex to help organize and translate my testing record into a structured validation report. I independently executed the team's code and replaced general AI-generated descriptions with the observed evidence: the \(2048\), \(4096\), and \(8192\) reference levels, the adjacent MAE reduction from \(0.0575\) to \(0.0407\), the finite-and-positive terminal-value checks, and the empirical slope of \(0.5111\). The numerical values and \texttt{PASS} results came from the program execution rather than from AI.\\[0.4em]

\textbf{What I would do differently / what I still need to understand:}\\
I would define the software-version requirements and evidence-storage structure before beginning the tests, then repeat the validation in an environment matching the exact dependency versions stated in the README. I would also add an automated test that runs on the designated team commit whenever later changes are introduced. I should have clarified earlier whether the testing role required formal visual inspection of every generated figure and whether exact README dependency versions were mandatory. My current checks establish the consistency of the main numerical results, but they do not fully evaluate figure presentation or every intermediate computational state.\\[0.6em]

\textbf{Signature:} \underline{\hspace{4cm}}\qquad
\textbf{Date:} \underline{\hspace{2.5cm}}
\end{formbox}

\newpage


% ============================================================
\section{Submission evidence matrix}
% ============================================================
\begin{longtable}{p{0.33\textwidth}p{0.58\textwidth}}
\toprule
\textbf{Course requirement} & \textbf{Evidence in this submission}\\
\midrule
SDE-specific methods and theory & Sections~\ref{sec:model}--\ref{sec:methods}:
exact GBM, EM, Milstein, strong/weak definitions, positivity, and coupled
non-affine EM.\\
Independent oracle & Exact transition, exact moments/distribution, analytic
mean bias, plus non-affine reference refinement.\\
Observed order for at least two methods & Table~\ref{tab:gbm-convergence} and
Figure~\ref{fig:gbm-convergence}, with slopes and fit uncertainty.\\
Monte Carlo uncertainty & Path-level 95\% confidence intervals in
Table~\ref{tab:nonaffine} and explicit formula~\eqref{eq:ci}.\\
Matched-accuracy cost & Section~\ref{sec:results}: common MAE target 0.85,
declared step-count cost, exclusions, and limitation.\\
Distribution and sanity checks & Figure~\ref{fig:gbm-distribution},
Table~\ref{tab:quantiles}, moment oracle, increment variance/correlation,
finiteness, positivity.\\
Non-affine nested validation & Equation~\eqref{eq:nested-sum},
Table~\ref{tab:nonaffine}, Figure~\ref{fig:nonaff-convergence}, and
2048-to-4096-to-8192 refinement and fixed-sample slope sensitivity
(Tables~\ref{tab:reference-refinement}--\ref{tab:reference-slopes}).\\
Reproducible code and figures & One command regenerates six fixed-name,
220-DPI figures, JSON, and console output; fixed seed; no absolute paths.\\
Interpretive visualisations & Every figure has units, a colorblind-safe palette,
and a caption stating both the evidence and its significance.\\
Scientific limitations & Section~\ref{sec:limitations} separates numerical
convergence, reference error, sampling uncertainty, cost, and model risk.\\
AI log and ICS & Appendices C and D, with audit/falsifier records and five
member-specific forms.\\
\bottomrule
\end{longtable}

\end{document}
